\subsection{A direct Fourier gate for the surviving index hypotheses}

\begin{proposition}[P-F34: a large cubic Fourier amplitude and the remaining average-polynomial gate]
\label{prop:cubicfourierindexgate}
For the ordinary cubic plethysm function $f_n$ of P-F31, let
$L$ be any quasipolynomial period and write
\[
 f_n(s)=q_r(s)\quad\hbox{when }s\equiv r\pmod L,
 \qquad
 B_n(x)=\frac1L\sum_{r=0}^{L-1}q_r(x).
 \tag{P-F34a}
\]
The polynomial $B_n$ is independent of the chosen period.
It is the contribution from the pole of $H_n(t)$ at $t=1$.
Its degree is
\[
 D=\binom{n+2}{3}-n^2=\binom n3.
\]
Put $M=2^{n-1}$ and
\[
 \mathcal C_n(\zeta)=\lim_{t\to\zeta}(1-t^M)H_n(t),
 \qquad A_n(\zeta)=\mathcal C_n(\zeta)/M
\]
for a primitive $M$th root $\zeta$. Then the following are
necessary conditions on the source's literal indices:
\[
 B_n\bigl(p(f_n)\bigr)\ge |A_n(\zeta)|,
 \qquad
 B_n\bigl(s(f_n)+1\bigr)\ge |A_n(\zeta)|.
 \tag{P-F34b}
\]
In the standard complex embedding
$\zeta=\exp(2\pi i/M)$, the amplitude satisfies, for every
$n\ge250$,
\[
                         |A_n(\zeta)|\ge2^{n^4/192}.
 \tag{P-F34c}
\]

Write $B_n(x)=\sum_{j=0}^{D}b_{n,j}x^j$ and
$U_n=\sum_j|b_{n,j}|>0$. Consequently
\[
 \max\{1,p(f_n)\},\quad s(f_n)+1
      \ \ge\ \left(\frac{2^{n^4/192}}{U_n}\right)^{1/D}
       \qquad(n\ge250).
 \tag{P-F34d}
\]
In particular, an additional estimate
\[
 U_n\le2^{C n^3\log_2(n+2)}
 \tag{P-F34e}
\]
with an absolute constant $C$ would imply, for all sufficiently
large $n$,
\[
 p(f_n)\ge\frac{2^{n/32}}{(n+2)^{6C}},\qquad
 s(f_n)+1\ge\frac{2^{n/32}}{(n+2)^{6C}}.
 \tag{P-F34f}
\]
Estimate (P-F34e) is not proved or assumed here. Thus the
proposition establishes a quantitative necessary test and
a sufficient remaining condition for an index refutation;
it does not itself refute narrow PH2 or SH.
\end{proposition}

\begin{proof}
We first give the general Fourier calculation, keeping the
normalization of a simple pole explicit. If $f$ is a
quasipolynomial of period $L$, with constituent polynomials
$q_r(x)$, define for each $L$th root of unity $\rho$
\[
 P_\rho(x)=\frac1L\sum_{r=0}^{L-1}\rho^r q_r(x).
 \tag{P-F34g}
\]
Finite Fourier inversion gives
$q_r(x)=\sum_{\rho^L=1}\rho^{-r}P_\rho(x)$, and hence
\[
 \sum_{s\ge0}f(s)t^s
   =\sum_{\rho^L=1}\sum_{s\ge0}P_\rho(s)(t/\rho)^s.
 \tag{P-F34h}
\]
For a nonzero polynomial $P$ of degree $a$, the generating
function $\sum_{s\ge0}P(s)z^s$ has its only pole at $z=1$,
of order $a+1$. One proof is to expand $P(s)$ in the
basis $\binom{s+j}{j}$, whose generating functions are
$(1-z)^{-j-1}$; the coefficient of the top basis member
is nonzero. Thus a simple pole of the left side at
$\rho=\zeta$ means precisely that $P_\zeta$ is a nonzero
constant. If $\zeta$ is primitive of order $M\mid L$,
its summand is $P_\zeta/(1-t/\zeta)$, and
\[
 \lim_{t\to\zeta}
       \frac{1-t^M}{1-t/\zeta}=M.
 \tag{P-F34i}
\]
For the family at hand, P-F31 therefore proves
\[
                   P_\zeta(x)=A_n(\zeta)
                         \quad\hbox{for every }x.
 \tag{P-F34j}
\]
There is no dependence on $x$ in this Fourier coefficient.

The coefficient for $\rho=1$ is $B_n$. Enlarging the
period only repeats the constituent polynomials, so this
average is intrinsic. Formula (P-F34h) also shows directly
that its generating function is the principal part of
$H_n$ at $t=1$. More explicitly, if that principal part is
\[
 \sum_{j=1}^{D+1}\frac{c_{n,j}}{(1-t)^j},
 \quad\hbox{then}\quad
 B_n(x)=\sum_{j=1}^{D+1}c_{n,j}
                        \binom{x+j-1}{j-1}.
 \tag{P-F34k}
\]
No exponential-period enumeration is needed to specify
this one polynomial once the Laurent coefficients at one
are known. No efficient procedure for finding them is
asserted.
The ring in P-F31 has dimension
$m-n^2+1$, with $m=\binom{n+2}{3}$, by the smooth tangent
calculation in P-F24. Its Hilbert series has a pole of
that order at one. This is the usual Hilbert dimension
statement for a finitely generated positively graded
domain, obtainable by taking a common Veronese and its
Hilbert polynomial. Consequently $B_n$ has degree
$D=m-n^2=\binom n3$.

For any real $x$ at which all $q_r(x)$ are nonnegative,
(P-F34g) and the triangle inequality imply
\[
 |A_n(\zeta)|
   =\left|\frac1L\sum_r\zeta^r q_r(x)\right|
   \le\frac1L\sum_r q_r(x)=B_n(x).
 \tag{P-F34l}
\]
After a positivity shift $p=p(f_n)$, every polynomial
$q_r(x+p)$ has nonnegative coefficients. Its constant
coefficient $q_r(p)$ is therefore nonnegative.
After a saturation shift $s=s(f_n)$, every nonzero
polynomial $q_r(x+s)$ is positive at every positive
integer $x$; at $x=1$ it is positive, while an identically
zero constituent has value zero. Applying (P-F34l) at
$p$ and $s+1$ proves (P-F34b).
These are the literal conventions of GCT VI,
Definition~1.2.2, printed p.~13, and Definition~1.2.4,
printed p.~14. In particular the saturation definition
requires positivity of each nonzero constituent at every
positive integer, not only at integers in its original
residue class.

We now prove the absolute-value estimate, rather than
trying to infer it from the cyclotomic valuation alone.
Put $g=\gcd(n,3)$, $\gamma=3/g$, and
$S_0=(1-(-2)^n)/3$. Choose
\[
 \zeta=e^{2\pi i/M},\qquad \xi=\zeta^{-g},\qquad
 A=\operatorname{diag}
        (\xi,\xi^{-2},\ldots,\xi^{(-2)^{n-1}}).
 \tag{P-F34m}
\]
The induced grading parameter is exactly $\zeta$:
\[
 \xi^{-\gamma S_0}=\zeta^{g\gamma S_0}
                 =\zeta^{1-(-2)^n}=\zeta.
\]
Here $g$ is odd and at most three, so $\xi$ is still
primitive of order $M$.
Let $w_i=(-2)^i$ for $0\le i\le n-2$, and let
$w_{n-1}=0$; these are integer representatives for the
chain weights modulo $M$.

The exact moving tangent sequence in P-F24 and P-F28
gives the following absolute-value identity:
\[
 |\mathcal C_n(\zeta)|
   =\frac{\displaystyle\prod_{i\ne j}
                     |1-\xi^{w_i-w_j}|}
          {\displaystyle\prod_{\substack{0\le i\le j\le k<n\\
                    w_i+w_j+w_k\not\equiv0\pmod M}}
                     |1-\xi^{w_i+w_j+w_k}|}.
 \tag{P-F34n}
\]
Indeed the moving cubic representation minus the moving
adjoint representation is the moving quotient tangent
representation, whose dual is the conormal space.
Determinants multiply in the resulting exact sequences.
Replacing a root-of-unity eigenvalue by its inverse
leaves $|1-\alpha|$ unchanged, so possible dual or
point-action conventions do not affect (P-F34n).
There are $n(n-1)$ numerator factors and $m-n$
denominator factors, since precisely the $n$ chain
cubic monomials are fixed.

For every nontrivial $M$th root $\alpha$,
\[
                         |1-\alpha|\ge4/M.
 \tag{P-F34o}
\]
To verify this, choose its angle with distance
$2\pi a/M$ from an integer multiple of $2\pi$,
where $1\le a\le M/2$. Then
$|1-\alpha|=2\sin(\pi a/M)\ge4a/M\ge4/M$,
using concavity of sine on $[0,\pi/2]$.
All numerator factors in (P-F34n) are nontrivial,
because the chain eigenvalues are distinct.

Let $q=\lfloor n/2\rfloor$, and use only cubic monomials
in the first $q$ variables, indexed $0,\ldots,q-1$.
There are
\[
 L_q=\binom{q+2}{3}-(q-1)
 \tag{P-F34p}
\]
moving monomials among them: exactly the $q-1$ consecutive
chain monomials are fixed. For each of these moving
monomials, its integer weight $b$ satisfies
$|b|\le3\cdot2^{q-1}$. Therefore
\[
 |1-\xi^b|=|1-e^{-2\pi igb/M}|
    \le \frac{2\pi g|b|}{M}
    \le \frac{9\pi\,2^q}{M}.
 \tag{P-F34q}
\]
Every other denominator factor is at most two.
Combining (P-F34n)--(P-F34q), and dividing by $M$
to pass from $\mathcal C_n$ to $A_n$, gives
\begin{align*}
 \log_2|A_n(\zeta)|
 &\ge L_q\bigl(n-q-\log_2(9\pi)\bigr)
       -n(n-1)(n-3)-m+1.
 \tag{P-F34r}
\end{align*}
This estimate is entirely in the specified complex
embedding.

For completeness, explicit elementary constants yield
(P-F34c). For $n\ge10$ one has
\[
 L_q\ge\frac{q^3}{6}\ge\frac{(n-1)^3}{48},\qquad
 n-q-\log_2(9\pi)\ge n/2-5\ge0,
\]
since $9\pi<32$. Substitution in (P-F34r) and expansion
of $m=n(n+1)(n+2)/6$ give
\[
 \log_2|A_n(\zeta)|
 \ge
 \frac{n^4-125n^3+369n^2-351n+106}{96}.
 \tag{P-F34s}
\]
For $n\ge250$, the first two terms in the numerator
are at least $n^4/2$, and
$369n^2-351n+106>0$. Thus the right side is at least
$n^4/192$. This proves (P-F34c).

Finally, for $x\ge0$,
\[
 |B_n(x)|\le U_n\max\{1,x\}^{D}.
\]
Together with (P-F34b)--(P-F34c), this proves
(P-F34d). If the additional bound (P-F34e) holds,
then for sufficiently large $n$ the quantity
$n^4/192-Cn^3\log_2(n+2)$ is positive. Since
$D=\binom n3\le n^3/6$, (P-F34d) gives
\[
 \log_2\max\{1,p(f_n)\},\quad\log_2(s(f_n)+1)
       \ge n/32-6C\log_2(n+2).
\]
The right hand side tends to infinity, so the maximum
with one can be removed for all sufficiently large
$n$. This proves (P-F34f), as a conditional consequence
only.
\end{proof}

\paragraph{The exact remaining estimate.}
The pole-at-one polynomial in (P-F34k) is a smaller
piece of data than the entire family of constituent
polynomials over its exponentially large period.
An estimate specific to the cubic invariant ring,
such as (P-F34e), would combine with the proved large
Fourier amplitude to refute polynomial index bounds.
No such estimate is supplied by P-F34. Its necessary
inequality is not an equivalence to PH2 or SH.

In particular $f_n(0)=1$, or bounds on finitely many
small-stretch coefficient values, do not bound $U_n$:
those values sample different constituents and do not
control their values away from their residue classes.
P-F33 supplies a concrete warning: it has actual large
primitive-root amplitudes and a growing average constant,
yet both strict indices are zero, even with normality,
rational singularities, Gorensteinness and negative
$a$-invariant. Thus no generic purity, normality or
Gorenstein argument can replace the remaining
ring-specific estimate. Large cyclotomic valuation
also cannot replace the absolute-value calculation
above, because the complex conjugate amplitudes may
have very different sizes.

The exact source for the index conventions is
\url{https://arxiv.org/pdf/0704.0229}, Definitions~1.2.2
and 1.2.4, printed pp.~13--14. The Fourier identity
and the quantitative bounds required here are proved
in full above; no historical priority is claimed.


\paragraph{A sharper asymptotic amplitude and a weaker sufficient estimate.}
The same product gives the more precise asymptotic formula
\[
 \log_2|A_n(e^{2\pi i/M})|=\frac{n^4}{24}+O(n^3).
 \tag{P-F34t}
\]
Here is a full proof, including control of the factors that are
not close to one. For $n\ge6$, set $K=n-6$ and first use only
monomials in variables with nonzero chain weights. For a monomial
whose largest index is $k$, the number of moving monomials is
\[
 N_0=1,\qquad N_k=\binom{k+2}{2}-1=\frac{k^2+3k}{2}
                      \quad(k\ge1).
\]
The subtraction removes precisely $x_{k-1}^2x_k$.
For any remaining monomial with this largest index, its integer
weight $b$ satisfies
\[
                       2^{k-2}\le |b|\le3\cdot2^k.
\]
To check the lower bound, if the largest index occurs two or
three times, the largest term dominates the other terms by at
least $3\cdot2^{k-1}$. If it occurs once and both smaller indices
are at most $k-2$, the triangle inequality gives $2^{k-1}$.
If one smaller index is $k-1$, the sum of these two terms is
$(-2)^k/2$; the other smaller term has magnitude at most
$2^{k-2}$, unless both smaller indices are $k-1$, which is
exactly the removed fixed monomial. The case $k=0$ is the
single cube and satisfies the displayed bounds directly.

For $k\le K$, $g|b|\le9\cdot2^{n-6}\le M/2$.
The same elementary sine bounds as above therefore give
\[
 \frac{2^k}{M}\le |1-\xi^b|
       \le\frac{18\pi\,2^k}{M}<\frac{64\,2^k}{M}.
\]
In particular the base-two logarithm of each such factor is
$k-n+\epsilon$ with $1\le\epsilon\le7$.
The main contribution to the negative logarithm of the cubic
moving determinant is consequently
\begin{align*}
 S_n&=\sum_{k=0}^{n-6}(n-k)N_k\\
    &=\frac{n^4+6n^3-181n^2+798n-720}{24}
      =\frac{n^4}{24}+O(n^3).
\end{align*}
The exact equality follows from the standard sums of the first
three powers of $k$ after substituting
$N_k=(k^2+3k)/2$ for $k\ge1$ and $N_0=1$.
The total logarithmic error in these bulk factors is at most
$7m=O(n^3)$.

There are at most $3n^2$ remaining cubic factors. Indeed a
monomial either contains the last variable of weight zero,
giving at most $\binom{n+1}{2}$ choices, or its largest index
is one of $n-5,n-4,n-3,n-2$, giving at most
$4\binom n2$ choices. Their sum is less than $3n^2$.
Each moving factor, regardless of angle, lies between $4/M$
and two, so its logarithm has absolute value at most $n$.
Thus all these exceptional factors together contribute
$O(n^3)$. The $n(n-1)$ adjoint factors obey the same bound
and also contribute $O(n^3)$. Finally division by $M$ changes
the logarithm by only $n-1$. This proves (P-F34t) with an
absolute implied constant; for example the estimates just
given imply absolute error at most $12n^3$ for $n\ge6$.

It follows that the sufficient remaining condition can be
weakened from (P-F34e) to
\[
 \limsup_{n\to\infty}\frac{\log_2 U_n}{n^4}<\frac1{24}.
 \tag{P-F34u}
\]
To verify the implication with its quantifiers, choose
$\delta>0$ such that
$\log_2 U_n\le(1/24-2\delta)n^4$ for every sufficiently
large $n$. Formula (P-F34t) then gives
$\log_2(|A_n|/U_n)\ge\delta n^4$ for all sufficiently
large $n$. The already proved index inequalities and
$D\le n^3/6$ imply
\[
                   p(f_n),\quad s(f_n)+1\ge2^{6\delta n}
\]
for all sufficiently large $n$; the maximum with one in
(P-F34d) disappears because this lower bound exceeds one.
Condition (P-F34u), like (P-F34e), remains unproved for the
actual cubic ring. P-F38 excludes deriving it from bounded
constraint entries, cubic-scale dimension, and a bounded
containing box alone. The sharper amplitude therefore
strengthens the precise remaining condition without settling
that condition or either surviving index hypothesis.
